% META: {"id":"1SPE-TRIGO-CO-006","chapitre":"1SPE-TRIGONOMETRIE","type_objet":"corrige","exercice_ref":"1SPE-TRIGO-EX-006","capacites_codes":["C1"],"sources_inspiration":[],"status":"approved","origin":"generated","status_history":["generated","audited","accepted","approved"],"release_acceptance":"RELEASE_OWNER_BATCH_ACCEPTANCE","acceptance_closure_digest":"sha256:9b3ccf9a81c5520fbb7e03b7b2d3e7bf2057a02834b6d908bf3be5f49f3b553f","acceptance_receipt_digest":"sha256:4fad86f0282e0fa4e0419cca1ad6379ae7af5a280c04a4a90880f90a1c044242"}
% BEGIN-VERIFY
% from sympy import *
% assert cos(pi/3) == Rational(1,2)
% assert sin(pi/3) == sqrt(3)/2
% assert cos(2*pi/3) == Rational(-1,2)
% assert sin(2*pi/3) == sqrt(3)/2
% assert cos(pi) == -1
% assert sin(pi) == 0
% END-VERIFY
\begin{corrige}{1SPE-TRIGO-EX-006}

\textbf{1.} L'hexagone regulier divise le cercle en $6$ arcs égaux. L'angle au centre est $\dfrac{2\pi}{6} = \dfrac{\pi}{3}$.

\textbf{2.} Les mesures associees sont : $A_1 \leftrightarrow \dfrac{\pi}{3}$, $A_2 \leftrightarrow \dfrac{2\pi}{3}$, $A_3 \leftrightarrow \pi$, $A_4 \leftrightarrow \dfrac{4\pi}{3}$, $A_5 \leftrightarrow \dfrac{5\pi}{3}$.

\textbf{3.} Coordonnées :
\begin{itemize}
  \item $A_0 = (1\,;\,0)$
  \item $A_1 = \left(\dfrac{1}{2}\,;\,\dfrac{\sqrt{3}}{2}\right)$
  \item $A_2 = \left(-\dfrac{1}{2}\,;\,\dfrac{\sqrt{3}}{2}\right)$
  \item $A_3 = (-1\,;\,0)$
  \item $A_4 = \left(-\dfrac{1}{2}\,;\,-\dfrac{\sqrt{3}}{2}\right)$
  \item $A_5 = \left(\dfrac{1}{2}\,;\,-\dfrac{\sqrt{3}}{2}\right)$
\end{itemize}

\textbf{4.} $A_3 = (\cos\pi\,;\,\sin\pi) = (-1\,;\,0)$. \checkmark

\end{corrige}
